By \(\mathrm{thm:matched\mbox{-}minimax\mbox{-}frontier}\), uniform consistency is equivalent to
\[
 \frac d{n\log(ed)}\longrightarrow0.                               \tag{24}
\]
We prove that (24) is equivalent to \(d=o(n\log(en))\).  Put \(t=d/n\) and \(L_n=\log(en)\).  If \(t/L_n\to0\), then on any subsequence with \(d\geq n\),
\[
 \log(ed)\geq L_n,
 \qquad
 \frac d{n\log(ed)}\leq\frac t{L_n}\longrightarrow0.
\]
On a subsequence with \(d<n\), if \(d\geq\sqrt n\), then \(\log(ed)\geq L_n/2\) for all sufficiently large \(n\), so
\[
 \frac d{n\log(ed)}\leq\frac2{L_n}\longrightarrow0;
\]
if \(d<\sqrt n\), that ratio is at most \(n^{-1/2}\).  Thus (24) follows.

Conversely, assume (24).  If \(t<1\), then \(t/L_n\leq1/L_n\to0\).  If \(t\geq1\), define
\[
 \eta_n=\frac{t}{L_n+\log t}.
\]
Because \(L_n+\log t=\log(ed)\), equation (24) says \(\eta_n\to0\).  Since \(\log t\leq t\), eventually \(\eta_n\log t\leq t/2\).  The identity
\[
 t=\eta_nL_n+\eta_n\log t
\]
then gives \(t/L_n\leq2\eta_n\to0\).  This proves the consistency equivalence.

Multiplying the matched frontier by \(n\) shows that parametric mean-squared risk is equivalent, up to constants, to boundedness above and away from zero of
\[
 n\min\left\{1,\frac d{n\log(ed)}\right\}.                         \tag{25}
\]
If \(d=O(1)\) and \(d\geq2\), then \(d/\log(ed)\) is bounded above and away from zero, so (25) is bounded above and away from zero for all sufficiently large \(n\).  Conversely, if \(d\to\infty\) along a subsequence, then either the minimum in (25) equals one, in which case (25) equals \(n\to\infty\), or it equals \(d/[n\log(ed)]\), in which case (25) equals
\[
 \frac d{\log(ed)}\longrightarrow\infty.
\]
Thus parametric mean-squared risk holds exactly for bounded alphabets.

Finally, \(\mathrm{prop:identification\mbox{-}and\mbox{-}extension}\) gives surjectivity of \(P\mapsto\operatorname{Obs}(P)\) from \(\mathcal P_{d,\epsilon}\) onto \(\mathcal D_{d,\epsilon}^{\mathrm{obs}}\), and \(\mathrm{prop:causal\mbox{-}optimal\mbox{-}value\mbox{-}corollary}\) gives equality of the causal and observed targets.  Hence the causal and observed minimax experiments have identical attainable worst-case risks, and both equivalences transfer to estimation of \(V^\star(P)\) over \(\mathcal P_{d,\epsilon}\).